\section{补充理论与证明}
\label{app:proofs}

\makeatletter
\@ifundefined{theHequation}{}{%
  \renewcommand*{\theHequation}{\thesection.\arabic{equation}}}
\@ifundefined{theHtheorem}{}{%
  \renewcommand*{\theHtheorem}{app.\thesection.\arabic{theorem}}}
\@ifundefined{theHproposition}{}{%
  \renewcommand*{\theHproposition}{app.\thesection.\arabic{proposition}}}
\@ifundefined{theHcorollary}{}{%
  \renewcommand*{\theHcorollary}{app.\thesection.\arabic{corollary}}}
\@ifundefined{theHlemma}{}{%
  \renewcommand*{\theHlemma}{app.\thesection.\arabic{lemma}}}
\@ifundefined{theHremark}{}{%
  \renewcommand*{\theHremark}{app.\thesection.\arabic{remark}}}
\makeatother

本附录说明 DELB 与既有下界的关系，陈述补充 Fano 对照，并证明正文中的全部
理论结果。熵表达式中的对数均以 \(2\) 为底，\(\ln\) 表示自然对数。

\subsection{与既有信息论下界的关系}
\label{subsec:prior_bound_scope}

DELB 是因果整数预测的专门化与综合，而不是新的基础最大熵定律。
Shannon 信源编码与率失真论证、Bayes 风险不等式、Ziv--Zakai 界和
Fano 类误差界处理不同的操作对象
\citep{shannon1948,shannon_rate_distortion,chen_bayes_risk,
jeong_ziv_zakai,feder_merhav}。Massey 的整数熵--方差关系及后续改进给出
格点修正 \citep{massey_integer_entropy,rioul_massey}，Fang 等人给出连续
条件熵预测论证 \citep{fang_variance_bounds}。近期工作进一步把连续论证用于时间
序列复杂度排序、用熵推导有限时预测空间的维度下界，并以认证熵值基准测试有限样本
熵估计器 \citep{ayers2025complexity,hauser2026dmd,calcagnile2026benchmark}。
表~\ref{tab:prior_bound_matrix} 区分这些组成和邻近目标与本文的四部分预测体系。

\begin{table}[!t]
\caption{Theorem-level boundary between established results and the present
DELB formulation. ``Present role'' states how a result is used; it is not an
assertion of absolute bibliographic priority.}
\label{tab:prior_bound_matrix}
\centering
\scriptsize
\setlength{\tabcolsep}{3pt}
\begin{tabularx}{\textwidth}{
    >{\RaggedRight\arraybackslash}p{0.15\textwidth}
    >{\RaggedRight\arraybackslash}p{0.17\textwidth}
    >{\RaggedRight\arraybackslash}p{0.25\textwidth}
    >{\RaggedRight\arraybackslash}X}
\toprule
\textbf{Result family} & \textbf{Target and loss} &
\textbf{Established result} & \textbf{Present role and non-claim}\\
\midrule
Rate--distortion and general Bayes-risk bounds
\citep{shannon_rate_distortion,chen_bayes_risk,jeong_ziv_zakai}
& General source/estimand and specified distortion or loss
& Operational coding tradeoffs and decision-theoretic or
hypothesis-testing lower bounds, including bounds applicable to discrete
estimands
& Establishes that discrete estimation-error lower bounds are not new as a
class; DELB is distinguished by its exact integer-entropy envelope and
causal prediction form\\

Fano and converse relations
\citep{cover_thomas,feder_merhav}
& Finite-alphabet target and Hamming error
& Relations between equivocation and minimum classification-error
probability
& Classification-loss counterpart only; neither Fano nor its inversion is
claimed as DELB novelty\\

Maximum entropy and information projection
\citep{csiszar_iprojection,cover_thomas}
& Distributions under linear moment constraints
& Exponential-family extremizers and KL/information-projection identities
for entropy deficits
& Supplies the general reason that the moment-matched entropy deficit is a
KL divergence; the identity is not new in isolation\\

Massey and later refinements
\citep{massey_integer_entropy,rioul_massey}
& Integer-valued variable and second moment or variance
& Discrete-Gaussian entropy extremizer and the closed-form
\(1/12\)-corrected entropy bound
& Supplies the exact lattice envelope and conservative closed-form
ingredient; neither is claimed as newly discovered\\

Fang et al.\ \citep{fang_variance_bounds}
& Real-valued process and squared prediction error
& Conditional differential-entropy bounds, residual independence, and
Gaussian innovation attainability
& Continuous prediction analogue and source of the entropy--error logic;
the contribution here is its exact integer-lattice specialization\\

Recent entropy-based prediction limits
\citep{ayers2025complexity,hauser2026dmd}
& Real-valued time-series complexity ranking, or the dimension of a
finite-time operator-prediction space
& Extensions of the continuous conditional-entropy argument and
measure-theoretic-entropy constraints on predictive representation dimension
& Direct 2025--2026 neighbors, but neither targets integer-valued squared
prediction error nor inverts the discrete maximum-entropy envelope\\

Recent entropy-estimator benchmarking
\citep{calcagnile2026benchmark}
& Finite-sample estimation of entropy or entropy rate
& Controlled comparison of several estimators against certified entropy
values for dynamical systems
& Supports the need to separate a population entropy estimand from estimator
behavior; it is not a prediction-error lower-bound theorem\\

Present DELB formulation
& Causal integer-valued prediction and squared error
& Exact inverse lattice envelope, predictor-dependent refinement, equality
conditions, and two-component prediction-slack diagnostic
& Claimed as a unified specialization and diagnostic synthesis, not a new
primitive entropy, KL, rate--distortion, Fano, or generic Bayes-risk
inequality\\
\bottomrule
\end{tabularx}
\end{table}


表~\ref{tab:core_bound_comparison} 给出用于说明正文方法定位的紧凑操作比较。

\begin{table}[H]
\caption{用于界定 DELB 理论贡献的主要下界框架比较。“精确”是指对离散
最大熵包络的精确反演，并不表示 DELB 对任意目标分布都等于 Bayes 风险。}
\label{tab:core_bound_comparison}
\centering
\scriptsize
\setlength{\tabcolsep}{3pt}
\begin{tabularx}{\textwidth}{
    >{\RaggedRight\arraybackslash}p{0.15\textwidth}
    >{\RaggedRight\arraybackslash}p{0.16\textwidth}
    >{\RaggedRight\arraybackslash}p{0.19\textwidth}
    >{\RaggedRight\arraybackslash}p{0.22\textwidth}
    >{\RaggedRight\arraybackslash}X}
\toprule
\textbf{框架} & \textbf{目标与损失} & \textbf{下界对象} &
\textbf{可达性与预测器作用} & \textbf{与本文 DELB 的关系}\\
\midrule
连续熵预测界 \citep{fang_variance_bounds}
& 实值目标；平方损失
& 通过高斯熵--方差关系把条件微分熵映射为误差界
& 等号要求创新具有相应高斯形式和独立性
& 是连续情形的对应结果，但不提供精确整数格点包络及格点修正\\

Massey 格点界 \citep{massey_integer_entropy,rioul_massey}
& 整数值变量；方差或二阶矩
& 离散最大熵以及带 \(1/12\) 修正的保守闭式界
& 给出格点极值分布，但本身不规定因果预测器、历史依赖强化或侧信息排序
& 提供既有格点组成；本文不把极值分布或闭式不等式声明为原创\\

Fano 不等式 \citep{cover_thomas,feder_merhav}
& 有限类别目标；Hamming 损失
& 把条件熵与分类错误概率联系起来
& 在正确反函数分支上约束最优分类决策
& 是损失特定的补充检验，不是 MSE 下界，也不属于 DELB 原创性主张\\

一般 Bayes 风险与率失真界
\citep{shannon_rate_distortion,chen_bayes_risk}
& 一般估计对象及给定损失或失真
& 通过决策或编码论证约束最小风险或可操作失真
& 在决策类上定义优化；紧致性取决于具体下界和统计实验
& 范围比 DELB 更广。DELB 是整数平方风险的显式信息论下界，而不是 Bayes
风险本身\\

本文 DELB 体系
& 因果整数值目标；平方损失
& 精确逆格点包络及预测器相关强化
& 给定预测器达到普适下界，当且仅当其误差为二阶矩匹配的离散高斯且与合法历史
独立；与最优风险相等还要求最优预测器存在
& 增加双松弛诊断、新增信息的 CMI 排序及有限样本估计--随机化流程；
其贡献是统一综合而非新的基础熵不等式\\
\bottomrule
\end{tabularx}
\end{table}


作为补充，表~\ref{tab:delb_theorem_package} 区分 DELB 体系回答的四类问题。

\begin{table}[!htbp]
\caption{DELB 四部分定理体系及各部分回答的不同问题。贡献在于针对因果整数
预测的联合组织；基础最大熵和信息投影结果仍属于既有理论。}
\label{tab:delb_theorem_package}
\centering
\footnotesize
\setlength{\tabcolsep}{4pt}
\begin{tabularx}{\textwidth}{
    >{\RaggedRight\arraybackslash}p{0.19\textwidth}
    >{\RaggedRight\arraybackslash}p{0.28\textwidth}
    >{\RaggedRight\arraybackslash}X}
\toprule
\textbf{组成} & \textbf{数学对象} & \textbf{回答的预测问题}\\
\midrule
精确逆格点底线
&
\(\mathcal H_{\mathbb Z}^{-1}
 [H(X_t\mid\mathcal F_{t-m})]\)
&
目标不确定性和整数格点对所有算法施加何种 MSE 下界？\\

预测器相关强化
&
\(\mathcal H_{\mathbb Z}^{-1}
 [H(X_t\mid\mathcal F_{t-m})+I(E_t;\mathcal F_{t-m})]\)
&
特定预测器的误差仍保留历史信息时，下界如何加强？\\

可达充要条件
&
\(P_E=q_D\) 且 \(I(E;\mathcal F)=0\)
&
预测器何时达到普适 DELB，而不仅仅是满足该不等式？\\

双松弛诊断
&
\(I(E;\mathcal F)+D_{\mathrm{KL},2}(P_E\Vert q_D)\)
&
不可达来自残余历史依赖、误差形状失配，还是二者兼有？\\
\bottomrule
\end{tabularx}
\end{table}


\begin{Remark}[贡献层次]
\label{rem:claim_hierarchy}
本文贡献在于为因果整数预测联合组织精确逆格点底线、预测器相关强化、可达充要条件
和双松弛分解，并把新增侧信息连接到可复刻的有限样本推断链。离散高斯极值分布、
KL 恒等式、格点修正和 Fano 不等式仍是继承组成。
\end{Remark}

\subsection{闭式、时间与向量扩展}
\label{subsec:closed_form_temporal}

\begin{Corollary}[闭式 DELB]
\label{cor:closed_form_delb}
定理~\ref{thm:exact_delb} 中每个预测器均满足
\begin{equation}
D_t^{(m)}
\geq
\left[
\frac{2^{2H(X_t\mid\mathcal F_{t-m})}}{2\pi e}
-\frac1{12}
\right]_+ .
\label{eq:discrete_mse_bound}
\end{equation}
\end{Corollary}

\(1/12\) 是均匀去量化产生的格点修正。该显式表达相对于精确逆格点 DELB
较保守，通常不可达。

\begin{Proposition}[预测提前期单调性]
\label{prop:general_horizon_monotonicity}
若 \(1\leq\ell<m\)，则
\begin{equation}
H(X_t\mid\mathcal F_{t-m})
\geq H(X_t\mid\mathcal F_{t-\ell}),
\label{eq:general_horizon_entropy}
\end{equation}
因而
\begin{equation}
\mathcal H_{\mathbb Z}^{-1}[H(X_t\mid\mathcal F_{t-m})]
\geq
\mathcal H_{\mathbb Z}^{-1}[H(X_t\mid\mathcal F_{t-\ell})].
\label{eq:general_horizon_delb}
\end{equation}
\end{Proposition}

\begin{Corollary}[平稳一步极限]
\label{cor:general_stationary_delb}
对严格平稳整数过程，若
\(\mathcal F_{t-1}=\sigma(X_0,\ldots,X_{t-1})\)，定义
\begin{equation}
\overline H(X)
=\lim_{t\to\infty}H(X_t\mid X_0,\ldots,X_{t-1}).
\label{eq:general_entropy_rate}
\end{equation}
则
\begin{equation}
\liminf_{t\to\infty}
\mathbb E[(X_t-\widehat X_t)^2]
\geq
\mathcal H_{\mathbb Z}^{-1}[\overline H(X)].
\label{eq:general_stationary_exact_bound}
\end{equation}
\end{Corollary}

\begin{Corollary}[整数向量闭式 DELB]
\label{cor:general_vector_delb}
对 \(\mathbf X_t,\widehat{\mathbf X}_t\in\mathbb Z^n\)，
\begin{equation}
\frac1n\mathbb E\|\mathbf X_t-\widehat{\mathbf X}_t\|_2^2
\geq
\left[
\frac{2^{(2/n)H(\mathbf X_t\mid\mathcal F_{t-m})}}{2\pi e}
-\frac1{12}
\right]_+ .
\label{eq:general_vector_delb}
\end{equation}
\end{Corollary}

上述结果分别使用均匀去量化、信息集嵌套和多元高斯最大熵不等式证明，
完整证明见附录~\ref{app:proofs}。


\subsection{分类损失下的 Fano 对照}
\label{subsec:fano_counterpart}

令固定量化器把整数目标映射为
\begin{equation}
Y_t=Q(X_t)\in\{1,\ldots,K\},
\label{eq:fano_quantized_target}
\end{equation}
令 \(\widehat Y_t\) 关于 \(\mathcal F_{t-m}\) 可测。记
\(P_e=\Pr(\widehat Y_t\neq Y_t)\)，并定义
\begin{equation}
\phi_K(p)=h_2(p)+p\log_2(K-1),
\qquad 0\leq p\leq1-\frac1K.
\label{eq:fano_phi}
\end{equation}

\begin{Proposition}[Fano 分类误差底线]
\label{prop:fano_error_floor}
每个上述分类器均满足
\begin{equation}
H(Y_t\mid\mathcal F_{t-m})
\leq h_2(P_e)+P_e\log_2(K-1),
\label{eq:fano_conditional_inequality}
\end{equation}
从而
\begin{equation}
P_e\geq
\phi_K^{-1}\!\left[H(Y_t\mid\mathcal F_{t-m})\right].
\label{eq:fano_general_floor}
\end{equation}
当 \(K=2\) 时，正确的单调分支给出
\begin{equation}
P_e\geq
h_2^{-1}\!\left[H(Y_t\mid\mathcal F_{t-m})\right],
\qquad h_2^{-1}:[0,1]\rightarrow[0,1/2].
\label{eq:fano_binary_floor}
\end{equation}
\end{Proposition}

对第~\ref{subsec:methods_information_sets} 节定义的基准和增强信息集，定义
\begin{equation}
P_{\mathrm F}^{0}
=\phi_K^{-1}[H(Y\mid\mathcal F^0)],
\qquad
P_{\mathrm F}^{B}
=\phi_K^{-1}[H(Y\mid\mathcal F^B)].
\label{eq:fano_mobility_floors}
\end{equation}

\begin{Corollary}[Fano 底线的侧信息排序]
\label{cor:mobility_fano_ordering}
熵下降
\begin{equation}
H(Y\mid\mathcal F^0)-H(Y\mid\mathcal F^B)
=I(Y;Q_B(\mathbf B^{(q)})\mid\mathcal F^0)\geq0
\label{eq:fano_mobility_cmi}
\end{equation}
意味着
\begin{equation}
P_{\mathrm F}^{B}\leq P_{\mathrm F}^{0}.
\label{eq:fano_mobility_floor_order}
\end{equation}
\end{Corollary}

这是分类损失下的对照，而不是计数 MSE DELB 的第二次推导。二者的目标、
损失、等号条件和数值尺度均不同。

\subsection{完整证明}

\subsection{条件整数平移}

\begin{Lemma}[可测整数平移下的熵不变性]
\label{lem:conditional_translation}
令 \(X\) 为整数值随机变量，\(\mathcal F\) 为 sigma 代数，
\(\widehat X\) 为整数值且关于 \(\mathcal F\) 可测。若
\(E=X-\widehat X\)，则
\begin{equation}
H(E\mid\mathcal{F})=H(X\mid\mathcal{F}).
\label{eq:conditional_translation}
\end{equation}
因此，
\begin{equation}
H(E)
=
H(X\mid\mathcal{F})+I(E;\mathcal{F}).
\label{eq:error_entropy_identity}
\end{equation}
\end{Lemma}

\begin{proof}
给定 \(\mathcal F\) 中信息的一次实现后，\(\widehat x\) 为固定值，且
\[
\Pr(E=k\mid\mathcal{F})
=
\Pr(X=k+\widehat{x}\mid\mathcal{F}).
\]
映射 \(k\mapsto k+\widehat x\) 是 \(\mathbb Z\) 上的双射，只会重新标记
条件概率质量函数的原子。Shannon 熵在这种重标记下保持不变。
对 \(\mathcal F\) 取平均即可得到式~\eqref{eq:conditional_translation}。
最后，由
\(I(E;\mathcal{F})=H(E)-H(E\mid\mathcal{F})\)
得到式~\eqref{eq:error_entropy_identity}。
\end{proof}

\subsection{整数最大熵分布}

\begin{Lemma}[离散高斯极值分布]
\label{lem:discrete_gaussian_extremizer}
对每个 \(D>0\)，定义~\ref{def:integer_entropy_envelope} 的熵包络由
式~\eqref{eq:general_discrete_gaussian} 的中心离散高斯分布取得，其中
唯一参数 \(\lambda>0\) 满足 \(D(\lambda)=D\)。其熵由
式~\eqref{eq:general_discrete_gaussian_entropy} 给出。当 \(D=0\) 时，
极值分布为零点处的点质量。
\end{Lemma}

\begin{proof}
记
\[
\Theta(\lambda)
=
\sum_{k\in\mathbb{Z}}\exp(-\lambda k^2),
\qquad
q_\lambda(k)
=
\frac{\exp(-\lambda k^2)}{\Theta(\lambda)}.
\]
对任意 \(\lambda>0\) 均可逐项求导，因此
\[
\mathbb{E}_{q_\lambda}[Z^2]
=
-\frac{\mathrm{d}}{\mathrm{d}\lambda}\ln\Theta(\lambda)
=D(\lambda).
\]
并且
\[
\frac{\mathrm{d}D(\lambda)}{\mathrm{d}\lambda}
=
-\operatorname{Var}_{q_\lambda}(Z^2)<0.
\]
当 \(\lambda\downarrow0\) 时 \(D(\lambda)\to\infty\)，当
\(\lambda\uparrow\infty\) 时 \(D(\lambda)\to0\)，故每个 \(D>0\)
对应唯一 \(\lambda\)。

令 \(p\) 为 \(\mathbb Z\) 上任意满足
\(\mathbb{E}_{p}[Z^2]\leq D\) 的概率质量函数，并选择使
\(D(\lambda)=D\) 的 \(\lambda\)。由相对熵非负性，
\begin{align}
0
&\leq
D_{\mathrm{KL}}(p\|q_\lambda)
=
\sum_{k\in\mathbb{Z}}p(k)
\log_2\frac{p(k)}{q_\lambda(k)}
\nonumber\\
&=
-H(p)
+
\frac{\lambda\mathbb{E}_{p}[Z^2]+\ln\Theta(\lambda)}{\ln2}.
\label{eq:kl_discrete_gaussian}
\end{align}
因此
\begin{equation}
H(p)
\leq
\frac{\lambda D+\ln\Theta(\lambda)}{\ln2}
=
H(q_\lambda).
\label{eq:integer_max_entropy}
\end{equation}
取等号要求
\[
D_{\mathrm{KL}}(p\|q_\lambda)=0,
\qquad
\mathbb{E}_{p}[Z^2]=D,
\]
从而 \(p=q_\lambda\)。
把 \(q_\lambda\) 直接代入
\(-\sum_kq_\lambda(k)\log_2q_\lambda(k)\) 即得到
式~\eqref{eq:general_discrete_gaussian_entropy}。当 \(D=0\) 时，
约束迫使 \(Z=0\) 几乎必然成立，故熵为零。
\end{proof}

\subsection{精确 DELB 的证明}

\begin{proof}[定理~\ref{thm:exact_delb} 的证明]
在引理~\ref{lem:conditional_translation} 中取
\(\mathcal F=\mathcal F_{t-m}\)，得到
\begin{equation}
H(E_t)
=
H(X_t\mid\mathcal{F}_{t-m})
+
I(E_t;\mathcal{F}_{t-m}).
\label{eq:appendix_error_entropy}
\end{equation}
因为 \(E_t\in\mathbb Z\) 且
\(\mathbb E[E_t^2]=D_t^{(m)}\)，熵包络定义给出
\[
H(E_t)\leq\mathcal{H}_{\mathbb{Z}}(D_t^{(m)}).
\]
广义逆单调不减，故
\[
D_t^{(m)}
\geq
\mathcal{H}_{\mathbb{Z}}^{-1}\!\left(H(E_t)\right).
\]
结合式~\eqref{eq:appendix_error_entropy} 得到
式~\eqref{eq:general_predictor_dependent_delb}。再利用互信息非负性并应用
相同单调逆，得到式~\eqref{eq:general_universal_delb}。

由引理~\ref{lem:discrete_gaussian_extremizer}，预测器相关不等式当且仅当
\(E_t\) 为二阶矩 \(D_t^{(m)}\) 的中心离散高斯时取等号，其中包括零点
退化分布。普适不等式取等号还要求
\(I(E_t;\mathcal F_{t-m})=0\)，即预测误差独立于可用信息。
\end{proof}

\subsection{精确可达性分解的证明}

\begin{proof}[命题~\ref{prop:attainability_decomposition} 的证明]
下面的计算把标准最大熵信息投影恒等式专门化为矩匹配离散高斯
\citep{csiszar_iprojection,cover_thomas}。为使预测特定等号条件和两个
松弛项自足，仍在此给出完整推导。

先设 \(D>0\)，将矩匹配离散高斯写为
\[
q_D(k)
=
\frac{\exp(-\lambda_D k^2)}{\Theta(\lambda_D)},
\qquad
\mathbb{E}_{q_D}[Z^2]=D.
\]
由于每个 \(k\in\mathbb Z\) 都有 \(q_D(k)>0\)，直接展开以 bit 为单位的
相对熵得
\begin{align}
D_{\mathrm{KL},2}(P_E\Vert q_D)
&=
\sum_{k\in\mathbb{Z}}P_E(k)
\log_2\frac{P_E(k)}{q_D(k)}
\nonumber\\
&=
-H(E)
+
\frac{\lambda_D\mathbb{E}[E^2]
      +\ln\Theta(\lambda_D)}{\ln2}
\nonumber\\
&=
\mathcal{H}_{\mathbb{Z}}(D)-H(E).
\label{eq:appendix_shape_slack}
\end{align}
最后一步使用 \(\mathbb E[E^2]=D\) 和
式~\eqref{eq:general_discrete_gaussian_entropy}。引理
\ref{lem:conditional_translation} 给出
\[
H(E)=H(X_t\mid\mathcal{F})+I(E;\mathcal{F}).
\]
代入式~\eqref{eq:appendix_shape_slack} 即得
式~\eqref{eq:predictor_attainability_decomposition}；两端再加
\(I(E;\mathcal F)\) 即得
式~\eqref{eq:universal_attainability_decomposition}。

当 \(D=0\) 时，\(E=0\) 几乎必然，故
\(q_0=P_E=\delta_0\)、\(H(E)=0\)、\(I(E;\mathcal F)=0\)，并由条件
平移有 \(H(X_t\mid\mathcal F)=0\)。两个等式均退化为 \(0=0\)。

相对熵与互信息均非负。因此，第一个松弛当且仅当 \(P_E=q_D\) 时为零；
普适松弛当且仅当 \(P_E=q_D\) 且 \(I(E;\mathcal F)=0\) 时为零。
对离散随机变量，零互信息等价于独立。
\end{proof}

\subsection{Fano 分类误差结果}

\begin{proof}[命题~\ref{prop:fano_error_floor} 的证明]
令 \(J=\mathbf 1\{\widehat Y_t\neq Y_t\}\)。由于 \(J\) 由
\((Y_t,\widehat Y_t)\) 确定，
\begin{align*}
H(Y_t\mid\widehat{Y}_t)
&=
H(J,Y_t\mid\widehat{Y}_t)\\
&=
H(J\mid\widehat{Y}_t)
+
H(Y_t\mid J,\widehat{Y}_t).
\end{align*}
第一项不超过 \(h_2(P_e)\)。当 \(J=0\) 时第二项为零；当 \(J=1\) 时，
至多剩余 \(K-1\) 个可能标签。因此
\[
H(Y_t\mid\widehat{Y}_t)
\leq
h_2(P_e)+P_e\log_2(K-1).
\]
由于 \(\widehat Y_t\) 关于 \(\mathcal F_{t-m}\) 可测，给定完整信息集
不会增加熵：
\[
H(Y_t\mid\mathcal{F}_{t-m})
\leq
H(Y_t\mid\widehat{Y}_t).
\]
由此得到式~\eqref{eq:fano_conditional_inequality}。

\(\phi_K\) 在 \([0,1-1/K]\) 上连续严格递增，取值从 \(0\) 到
\(\log_2K\)。若 \(P_e\leq1-1/K\)，单调求逆得到
式~\eqref{eq:fano_general_floor}；若 \(P_e>1-1/K\)，因逆函数取值仅在
\([0,1-1/K]\)，该式自然成立。当 \(K=2\) 时
\(\log_2(K-1)=0\)，所以 \([0,1/2]\) 上 \(\phi_2=h_2\)，从而得到
式~\eqref{eq:fano_binary_floor}。
\end{proof}

\begin{proof}[推论~\ref{cor:mobility_fano_ordering} 的证明]
根据 \(\mathcal F^B_{c,i,t}\) 的定义，条件互信息恒等式给出
\[
H(Y_{c,i,t}\mid\mathcal{F}^{0}_{c,i,t})
-
H(Y_{c,i,t}\mid\mathcal{F}^{B}_{c,i,t})
=
I\!\left(
Y_{c,i,t};
Q_B(\mathbf{B}^{(q)}_{c,i,t})
\mid\mathcal{F}^{0}_{c,i,t}
\right)
\geq0.
\]
\(\phi_K^{-1}\) 在递增分支上单调不减。对两个条件熵应用该函数，得到
式~\eqref{eq:fano_mobility_floor_order}。命题
\ref{prop:fano_error_floor} 分别适用于关于任一信息集可测的分类器。
\end{proof}

\subsection{均匀去量化与闭式界}

\begin{proof}[推论~\ref{cor:closed_form_delb} 的证明]
令 \(U\sim\operatorname{Unif}[-1/2,1/2)\) 独立于 \(E_t\)，并令
\(Y=E_t+U\)。在区间 \([k-1/2,k+1/2)\) 上，\(Y\) 的密度等于
\(\Pr(E_t=k)\)。各区间互不相交且长度为 1，故
\begin{equation}
h(Y)=H(E_t).
\label{eq:dither_entropy_identity}
\end{equation}
此外，
\begin{equation}
\operatorname{Var}(Y)
=
\operatorname{Var}(E_t)+\frac{1}{12}
\leq
\mathbb{E}[E_t^2]+\frac{1}{12}
=
D_t^{(m)}+\frac{1}{12}.
\label{eq:dither_variance}
\end{equation}
高斯分布在固定方差下最大化微分熵，因此
\begin{equation}
H(E_t)
=h(Y)
\leq
\frac{1}{2}\log_2\!\left[
2\pi e\left(D_t^{(m)}+\frac{1}{12}\right)
\right].
\label{eq:dither_gaussian_bound}
\end{equation}
整理得
\[
D_t^{(m)}
\geq
\frac{2^{2H(E_t)}}{2\pi e}-\frac{1}{12}.
\]
式~\eqref{eq:appendix_error_entropy} 同时给出预测器相关版本，并在丢弃
非负互信息项后得到普适版本。最后使用 \(D_t^{(m)}\geq0\) 并取正部，
即可得到式~\eqref{eq:discrete_mse_bound}。
\end{proof}

\subsection{时间与向量推论}

\begin{proof}[命题~\ref{prop:general_horizon_monotonicity} 的证明]
对 \(1\leq\ell<m\)，因果嵌套给出
\(\mathcal F_{t-m}\subseteq\mathcal F_{t-\ell}\)。因此
\[
H(X_t\mid\mathcal{F}_{t-m})
-
H(X_t\mid\mathcal{F}_{t-\ell})
=
I(X_t;\mathcal{F}_{t-\ell}\mid\mathcal{F}_{t-m})
\geq0,
\]
即式~\eqref{eq:general_horizon_entropy}。
\(\mathcal H_{\mathbb Z}^{-1}\) 的单调性给出
式~\eqref{eq:general_horizon_delb}。
\end{proof}

\begin{proof}[推论~\ref{cor:general_stationary_delb} 的证明]
由平稳性，
\[
H(X_t\mid X_0,\ldots,X_{t-1})
=
H(X_0\mid X_{-t},\ldots,X_{-1}).
\]
右端随 \(t\) 单调不增并收敛到熵率 \(\overline H(X)\)。
对每个 \(t\) 应用式~\eqref{eq:general_universal_delb}，再使用由引理
\ref{lem:discrete_gaussian_extremizer} 的离散高斯参数化所保证的逆熵
包络连续性与单调性，即得式~\eqref{eq:general_stationary_exact_bound}。
\end{proof}

\begin{proof}[推论~\ref{cor:general_vector_delb} 的证明]
令 \(\mathbf E_t=\mathbf X_t-\widehat{\mathbf X}_t\)，令
\(\mathbf U\) 独立且在 \([-1/2,1/2)^n\) 上均匀分布。单位立方体
\(\mathbf k+[-1/2,1/2)^n\) 互不相交，所以
\[
h(\mathbf{E}_t+\mathbf{U})=H(\mathbf{E}_t).
\]
令 \(\boldsymbol\Sigma_E\) 为 \(\mathbf E_t\) 的协方差矩阵。多元高斯
最大熵不等式和特征值算术--几何平均不等式给出
\begin{align}
H(\mathbf{E}_t)
&\leq
\frac{1}{2}\log_2\!\left[
(2\pi e)^n
\det\!\left(\boldsymbol{\Sigma}_E+\frac{1}{12}\mathbf{I}_n\right)
\right]
\nonumber\\
&\leq
\frac{n}{2}\log_2\!\left[
2\pi e
\left(
\frac{\operatorname{tr}(\boldsymbol{\Sigma}_E)}{n}
+
\frac{1}{12}
\right)
\right]
\nonumber\\
&\leq
\frac{n}{2}\log_2\!\left[
2\pi e
\left(
\frac{\mathbb{E}\|\mathbf{E}_t\|_2^2}{n}
+
\frac{1}{12}
\right)
\right].
\label{eq:vector_dither_bound}
\end{align}
条件整数向量平移是双射，因此
\[
H(\mathbf{E}_t)
\geq
H(\mathbf{E}_t\mid\mathcal{F}_{t-m})
=
H(\mathbf{X}_t\mid\mathcal{F}_{t-m}).
\]
代入式~\eqref{eq:vector_dither_bound}、整理并利用 MSE 非负性，即得
式~\eqref{eq:general_vector_delb}。必须使用联合条件熵；用边际熵之和替代
一般会丢失空间依赖。
\end{proof}

\subsection{犯罪计数与自行车移动的专门化}

犯罪计数结果由在定理~\ref{thm:exact_delb} 中直接代入
\(X_t=C_{c,i,t}\) 和
\(\mathcal F_{t-m}=\mathcal F^0_{c,i,t}\) 得到。这证明
式~\eqref{eq:local_bound_baseline} 与
\eqref{eq:crime_baseline_exact_bound}；推论~\ref{cor:closed_form_delb}
证明式~\eqref{eq:crime_baseline_closed_bound}。若目标被截顶，确定性数据
处理给出式~\eqref{eq:top_coded_entropy_order}。因为
\(\mathcal H_{\mathbb Z}^{-1}\) 单调不减，代入更小的截顶熵会产生不大于
原始界的下界，因此对原始计数 MSE 仍是保守有效的。

\begin{proof}[命题~\ref{prop:bicycle_information_value} 的证明]
由式~\eqref{eq:bike_information}，
\[
\mathcal{F}^{B}_{c,i,t}
=
\mathcal{F}^{0}_{c,i,t}
\vee
\sigma\!\left\{Q_B(\mathbf{B}^{(q)}_{c,i,t})\right\}.
\]
条件互信息定义给出
\begin{align*}
I\!\left(
C_{c,i,t};
Q_B(\mathbf{B}^{(q)}_{c,i,t})
\mid\mathcal{F}^{0}_{c,i,t}
\right)
&=
H(C_{c,i,t}\mid\mathcal{F}^{0}_{c,i,t})\\
&\quad-
H(C_{c,i,t}\mid\mathcal{F}^{B}_{c,i,t}),
\end{align*}
这就是非负的式~\eqref{eq:bicycle_information_gain}。应用单调不减函数
\(\mathcal H_{\mathbb Z}^{-1}\)，得到
式~\eqref{eq:bicycle_exact_bound_order}。根据
式~\eqref{eq:discrete_entropy_power}，
\[
\frac{N_d(C\mid\mathcal{F}^{B})}
{N_d(C\mid\mathcal{F}^{0})}
=
2^{\,2[H(C\mid\mathcal{F}^{B})-H(C\mid\mathcal{F}^{0})]}
=
2^{-2\Delta H^B},
\]
从而得到式~\eqref{eq:cmi_entropy_power_ratio}。

最后，给定 \(\mathcal F^0_{c,i,t}\) 后，
\(Q_B(\mathbf B^{(q)}_{c,i,t})\) 是
\(\mathbf B^{(q)}_{c,i,t}\) 的确定性函数。因此条件 Markov 链
\[
C_{c,i,t}
\longrightarrow
\mathbf{B}^{(q)}_{c,i,t}
\longrightarrow
Q_B(\mathbf{B}^{(q)}_{c,i,t})
\]
满足条件数据处理不等式，从而证明
式~\eqref{eq:bicycle_data_processing}。
\end{proof}
